\documentclass[11pt, letterpaper]{article}
% ----- PAQUETES -----
\usepackage[utf8]{inputenc}
\usepackage[T1]{fontenc}
\usepackage[spanish]{babel}
\usepackage{graphicx}
\usepackage[margin=2.3cm]{geometry}
\usepackage{xcolor}
\usepackage{titlesec}
\usepackage{booktabs}
\usepackage{enumitem}
\usepackage{xurl}
\usepackage{hyperref}
\usepackage{fancyhdr}
\usepackage{amsmath}
\usepackage{amssymb}
\usepackage{array}
\usepackage{tikz}
\usetikzlibrary{arrows.meta, babel, calc, angles}

% ----- COLORES -----
\definecolor{UnalGreen}{RGB}{0,153,51}
\definecolor{DarkGray}{RGB}{60,60,60}
\definecolor{StripA}{RGB}{200,60,60}
\definecolor{StripB}{RGB}{60,90,180}

% ----- SECCIONES -----
\titleformat{\section}
{\normalfont\large\bfseries\color{UnalGreen}}
{\thesection}{1em}{}[{\titlerule[0.8pt]}]
\titleformat{\subsection}
{\normalfont\normalsize\bfseries\color{DarkGray}}
{\thesubsection}{1em}{}
\titlespacing*{\section}{0pt}{8pt}{4pt}
\titlespacing*{\subsection}{0pt}{6pt}{2pt}

% ----- ENLACES -----
\hypersetup{colorlinks=true, linkcolor=UnalGreen, urlcolor=UnalGreen}

% ----- ENCABEZADO -----
\pagestyle{fancy}
\fancyhf{}
\lhead{\footnotesize Mecánica y Ondas para Biociencias (1000023) -- Laboratorio}
\rhead{\footnotesize Semestre 2026-2}
\cfoot{\thepage}

% ----- LISTAS COMPACTAS -----
\setlist[itemize]{leftmargin=*, itemsep=2pt, topsep=2pt, parsep=0pt}
\setlist[enumerate]{leftmargin=*, itemsep=3pt, topsep=3pt, parsep=0pt}

% ----- CAJA DE NOTA -----
\newcommand{\notabox}[1]{%
	\begin{center}
		\fcolorbox{UnalGreen}{UnalGreen!8}{%
			\parbox{0.94\textwidth}{\small #1}%
		}
	\end{center}
}

% ----- LEYENDA DE FIGURAS -----
\newcounter{figg}
\newcommand{\leyenda}[2]{%
	\par\vspace{0.1cm}\refstepcounter{figg}\label{#1}%
	{\centering\small\textbf{Figura~\thefigg.} #2\par}\vspace{0.2cm}}

% =========================================================
% ESTILOS Y MACROS TikZ
% =========================================================
\tikzset{
	fuerza/.style={-{Latex[length=2.4mm]}, very thick, UnalGreen!70!black},
	peso/.style={-{Latex[length=2.4mm]}, very thick, StripA},
	reac/.style={-{Latex[length=2.4mm]}, very thick, StripB},
	dim/.style={{Latex[length=1.6mm]}-{Latex[length=1.6mm]}, thin, DarkGray},
}

% Barra dentada: 11 orificios, longitud 11 (unidades de dibujo).
% #1 = x del extremo izquierdo, #2 = y del eje central
\newcommand{\barra}[2]{%
	\draw[thick, fill=orange!20] (#1,#2-0.25) rectangle (#1+11,#2+0.25);
	\foreach \hi in {0,...,10}{\draw[thin, fill=white] (#1+0.5+\hi,#2) circle (0.09);}
}

% Pesa colgada de la barra: #1 = x, #2 = y de la cara inferior de la barra, #3 = etiqueta
\newcommand{\pesa}[3]{%
	\draw[thick] (#1,#2) -- (#1,#2-0.4);
	\draw[thick, fill=StripA!25] (#1-0.4,#2-1.1) rectangle (#1+0.4,#2-0.4);
	\node[font=\scriptsize] at (#1,#2-0.75) {#3};
}

% Apoyo (fulcro): triangulo bajo la barra y perno en el orificio #1
\newcommand{\fulcro}[1]{%
	\fill[gray!55] (#1,-0.25) -- (#1-0.5,-1.0) -- (#1+0.5,-1.0) -- cycle;
	\draw[thick] (#1-0.8,-1.0) -- (#1+0.8,-1.0);
	\draw[thick, fill=white] (#1,0) circle (0.11);
}

% Dinamometro que cuelga de un soporte superior (y = #2) y sujeta la barra en x = #1
\newcommand{\dinaarriba}[2]{%
	\draw[line width=2pt, gray!70] (#1-0.7,#2) -- (#1+0.7,#2);
	\draw[thick] (#1,#2) -- (#1,#2-0.2);
	\draw[thick, fill=white] (#1-0.22,#2-1.2) rectangle (#1+0.22,#2-0.2);
	\foreach \dy in {0.4,0.6,0.8,1.0}{\draw[thin] (#1-0.22,#2-\dy) -- (#1-0.08,#2-\dy);}
	\draw[thick] (#1,#2-1.2) -- (#1,0.25);
	\fill (#1,0) circle (0.05);
}

% Dinamometro que tira hacia abajo, anclado en y = #2 (negativo)
\newcommand{\dinaabajo}[2]{%
	\draw[line width=2pt, gray!70] (#1-0.7,#2) -- (#1+0.7,#2);
	\draw[thick] (#1,-0.25) -- (#1,-0.8);
	\draw[thick, fill=white] (#1-0.22,-1.9) rectangle (#1+0.22,-0.8);
	\foreach \dy in {1.1,1.3,1.5,1.7}{\draw[thin] (#1-0.22,-\dy) -- (#1-0.08,-\dy);}
	\draw[thick] (#1,-1.9) -- (#1,#2);
	\fill (#1,0) circle (0.05);
}

\begin{document}
	
	% =========================================================
	% ENCABEZADO
	% =========================================================
	\begin{center}
		\vspace*{-0.8cm}
		\IfFileExists{imagenes/universidad-nacional-de-colombia-sede-bogota-seeklogo.png}{%
			\includegraphics[width=0.15\textwidth]
			{imagenes/universidad-nacional-de-colombia-sede-bogota-seeklogo.png}\\[0.1cm]
		}{}
		
		{\Large\bfseries Práctica 5: Equilibrio y Momentos}\\[0.1cm]
		{\large\bfseries Torque, palancas y equilibrio de una barra}\\[0.1cm]
		{\normalsize Mecánica y Ondas para Biociencias (1000023)}
	\end{center}
	
	\vspace{0.2cm}
	\noindent
	\begin{tabular}{@{}ll@{}}
		Modalidad: & Grupal (3 a 4 integrantes) \\
		Entregable: & Informe 1 (formato en la Sección~\ref{sec:informe}) \\
		Fecha de sesión: & 20 de octubre de 2026 \\
		Primera versión del informe: & 3 de noviembre de 2026 \\
	\end{tabular}
	
	\notabox{\textbf{Antes de venir al laboratorio:} lean la descripción del
		experimento.}
	
	% =========================================================
	\section{Objetivos}
	% =========================================================
	\begin{itemize}
		\item Reconocer situaciones de equilibrio mecánico: cuando un objeto no
		se desplaza y cuando no gira.
		\item Usar el torque para \emph{predecir} la lectura de un dinamómetro.
		\item Comparar la predicción con la medida y clasificar el montaje como
		palanca de primer, segundo o tercer género.
		\item Distinguir equilibrio estable, indiferente e inestable.
	\end{itemize}
	
	% =========================================================
	\section{Materiales}
	% =========================================================
	\begin{itemize}
		\item Soporte universal.
		\item Barra dentada (barra con orificios equiespaciados).
		\item Juego de masas con portapesas.
		\item Dinamómetro.
		\item Balanza.
		\item Regla o cinta métrica.
	\end{itemize}
	
	% =========================================================
	\section{Descripción del experimento}
	% =========================================================
	
	Una puerta es más fácil de abrir si se empuja lejos de las bisagras que si
	se empuja cerca de ellas. Lo que hace girar un objeto no es solo la fuerza,
	sino la fuerza multiplicada por la distancia al punto de giro: el
	\textbf{torque}. Un objeto está en \textbf{equilibrio} cuando ni se
	desplaza ni gira.
	
	\vspace{0.3cm}
	\begin{center}
		\begin{tikzpicture}[scale=0.9, every node/.style={font=\small}]
			% ---------- (a) fuerza perpendicular ----------
			\begin{scope}
				\draw[line width=4pt, gray!60, line cap=round] (0,0) -- (4.2,0);
				\draw[thick, fill=white] (0,0) circle (0.14);
				\node[below left] at (0,0) {punto de giro};
				\draw[fuerza] (4.2,0) -- (4.2,-1.3) node[right] {$F$};
				\draw[dim] (0,0.5) -- (4.2,0.5) node[midway, above] {$d$};
				\draw[->, thick, StripB] (215:0.65) arc (215:145:0.65);
				\node[align=center] at (2.1,-2.1) {(a) La fuerza empuja\\ perpendicular a la barra};
			\end{scope}
			% ---------- (b) fuerza oblicua ----------
			\begin{scope}[xshift=7.4cm]
				\coordinate (O) at (0,0);
				\coordinate (P) at (4.2,0);
				\coordinate (Q) at ($(P)+(1.2,1.5)$);
				\coordinate (H) at ($(P)!(O)!(Q)$);
				\draw[line width=4pt, gray!60, line cap=round] (O) -- (P);
				\draw[thick, fill=white] (O) circle (0.14);
				\draw[dashed, gray] ($(P)!-1.45!(Q)$) -- ($(P)!1.3!(Q)$);
				\draw[fuerza] (P) -- (Q) node[above right] {$F$};
				\draw[very thick, StripA] (O) -- (H);
				\draw pic[draw, angle radius=2.2mm] {right angle=O--H--P};
				\node[StripA, below left] at ($(O)!0.5!(H)$) {$d$};
				\node[align=center] at (2.4,-2.9) {(b) La fuerza empuja\\ en diagonal};
			\end{scope}
		\end{tikzpicture}
	\end{center}
	\leyenda{fig:torque}{El torque es $\tau = F\times d$. La distancia $d$ (el
		\emph{brazo}) va del punto de giro a la línea sobre la que actúa la fuerza
		(línea punteada), medida \emph{en perpendicular} (línea roja). En (a) $d$ es
		toda la longitud de la barra; en (b) es más corta, y por eso girar cuesta más.}
	
	\vspace{0.1cm}
	\noindent
	El experimento se hace con la \textbf{barra dentada}: una barra con
	orificios equiespaciados, que sirven como puntos de apoyo y como puntos para
	colgar masas o el dinamómetro.
	
	\vspace{0.2cm}
	\begin{center}
		\begin{tikzpicture}[scale=0.85, every node/.style={font=\small}]
			\barra{0}{0}
			% numeracion de orificios
			\foreach \hi in {0,...,10}{
				\pgfmathtruncatemacro{\num}{\hi+1}
				\node[font=\scriptsize, gray!70!black] at (0.5+\hi,0.55) {\num};
			}
			% centro de gravedad
			\draw[thick, StripA] (5.5,-0.25) -- (5.5,0.25);
			\draw[peso] (5.5,-0.25) -- (5.5,-1.5) node[right] {$P_b$};
			\node[StripA, align=left, font=\scriptsize] at (5.5,-2.0) {peso de la barra, aplicado en el centro de gravedad (CG)};
			% separacion entre orificios
			\draw[dim] (0.5,-0.6) -- (1.5,-0.6);
			\node[font=\scriptsize] at (1.0,-0.95) {$a$};
		\end{tikzpicture}
	\end{center}
	\leyenda{fig:barra}{Barra dentada (ilustrativa: el número de orificios y su
		separación $a$ pueden ser distintos en su barra). El peso de la barra $P_b$
		actúa en su \emph{centro de gravedad}, que es el punto donde se puede pensar
		que está concentrado todo el peso.}
	
	% =========================================================
	\section{Teoría}
	% =========================================================
	
	\subsection*{Torque}
	
	El torque de una fuerza es
	\begin{equation}
		\tau = F\times d
		\label{eq:torque}
	\end{equation}
	
	\noindent
	donde $F$ es la fuerza y $d$ el brazo (Figura~\ref{fig:torque}). Sus
	unidades son N$\cdot$m (aquí usaremos N$\cdot$cm). Un torque puede hacer
	girar la barra en sentido \textbf{horario} (como las manecillas del reloj) o
	\textbf{antihorario}.
	
	\subsection*{Condiciones de equilibrio}
	
	Una barra queda quieta si se cumplen \textbf{las dos} condiciones:
	\begin{enumerate}
		\item \emph{No se desplaza:} la suma de las fuerzas hacia arriba es igual
		a la suma de las fuerzas hacia abajo.
		\item \emph{No gira:} respecto a cualquier punto, el torque total
		horario es igual al torque total antihorario.
	\end{enumerate}
	
	\noindent
	La condición de giro se puede escribir respecto a \emph{cualquier} punto.
	Conviene escoger el apoyo: la fuerza que hace el apoyo pasa por ese punto,
	tiene brazo cero y no produce torque, así que desaparece de la ecuación.
	
	\subsection*{Palancas}
	
	Una palanca es una barra que gira alrededor de un apoyo. Hay tres
	protagonistas: la \textbf{carga} $P$ (lo que se quiere sostener o mover), el
	\textbf{esfuerzo} $F$ (la fuerza que se aplica; en el laboratorio, la del
	dinamómetro) y el \textbf{apoyo}, que también hace una fuerza $N$. El género
	depende de cuál queda en medio:
	
	\begin{center}
		\begin{tikzpicture}[scale=0.62, every node/.style={font=\small}]
			\tikzset{barrita/.style={line width=3.5pt, orange!60!black, line cap=round}}
			\tikzset{apoyo/.style={fill=gray!55, draw=black, thick}}
			% ---------- 1er genero ----------
			\begin{scope}
				\draw[barrita] (0,0) -- (8,0);
				\draw[apoyo] (4,0) -- (3.5,-0.8) -- (4.5,-0.8) -- cycle;
				\draw[peso] (0.8,0) -- (0.8,-1.5) node[below] {$P$};
				\draw[fuerza] (7.2,0) -- (7.2,-1.5) node[below] {$F$};
				\draw[reac] (4,-0.8) -- (4,1.3) node[above] {$N$};
				\node at (4,-2.6) {1.er género};
			\end{scope}
			% ---------- 2do genero ----------
			\begin{scope}[xshift=9.5cm]
				\draw[barrita] (0,0) -- (8,0);
				\draw[apoyo] (0,0) -- (-0.5,-0.8) -- (0.5,-0.8) -- cycle;
				\draw[peso] (3,0) -- (3,-1.5) node[below] {$P$};
				\draw[fuerza] (7.2,0) -- (7.2,1.5) node[above] {$F$};
				\draw[reac] (0,-0.8) -- (0,1.3) node[above] {$N$};
				\node at (4,-2.6) {2.º género};
			\end{scope}
			% ---------- 3er genero ----------
			\begin{scope}[xshift=19cm]
				\draw[barrita] (0,0) -- (8,0);
				\draw[apoyo] (0,0) -- (-0.5,-0.8) -- (0.5,-0.8) -- cycle;
				\draw[peso] (7.2,0) -- (7.2,-1.5) node[below] {$P$};
				\draw[fuerza] (3,0) -- (3,1.5) node[above] {$F$};
				\draw[reac] (0,-0.8) -- (0,1.3) node[above] {$N$};
				\node at (4,-2.6) {3.er género};
			\end{scope}
		\end{tikzpicture}
	\end{center}
	\leyenda{fig:generos}{Los tres géneros de palanca: el apoyo queda en medio
		(1.º), la carga queda en medio (2.º) o el esfuerzo queda en medio (3.º). En
		el primero, el esfuerzo empuja hacia abajo, igual que la carga, pero del otro
		lado del apoyo; por eso en el laboratorio el dinamómetro tirará \emph{hacia
			abajo}. En este dibujo se desprecia el peso de la barra; en el laboratorio
		\textbf{no} se puede despreciar.}
	
	\noindent
	Sin contar el peso de la barra, la condición de giro respecto al apoyo da
	``torque del esfuerzo = torque de la carga'':
	\[
	F\times d_F = P\times d_P
	\]
	
	\noindent
	Si el esfuerzo está \emph{más lejos} del apoyo que la carga ($d_F > d_P$),
	basta una fuerza menor que el peso (2.º género). Si está \emph{más cerca}
	($d_F < d_P$), hace falta una fuerza mayor que el peso (3.º género).
	
	\subsection*{La fórmula de la práctica}
	
	En el laboratorio la carga es una masa $m$ colgada de la barra, y la barra
	tiene su propio peso $P_b$. Con el apoyo y el dinamómetro como en la
	Figura~\ref{fig:palancas}, los pesos de la masa y de la barra quedan del
	mismo lado del apoyo y el dinamómetro del otro. Entonces, el torque del
	dinamómetro iguala a la suma de los dos torques:
	\begin{equation}
		F_{\text{pred}}\times d_F = P_m\times d_m + P_b\times d_b
		\quad\Longrightarrow\quad
		F_{\text{pred}} = \frac{P_m\,d_m + P_b\,d_b}{d_F}
		\label{eq:Fpred}
	\end{equation}
	
	\noindent
	donde $P_m = m\,g$ es el peso de la masa, $P_b = m_b\,g$ el de la barra,
	$d_m$ y $d_F$ las distancias del apoyo a la masa y al dinamómetro, y $d_b$ la
	distancia del apoyo al centro de gravedad de la barra. Si el apoyo está
	exactamente en el centro de gravedad, $d_b = 0$ y el peso de la barra no
	produce torque.
	
	\notabox{Usen $m$ en \textbf{kg} y $g = 9.78$ m/s$^2$, de modo que $P$ queda
		en N. Las distancias en cm: así $F_{\text{pred}}$ sale directamente en N.}
	
	\subsection*{Tipos de equilibrio}
	
	Si a un cuerpo en equilibrio se le da una pequeña perturbación, ocurre una
	de tres cosas, según cómo cambie la altura de su centro de gravedad:
	
	\vspace{0.1cm}
	\begin{center}
		\begin{tikzpicture}[scale=0.9, every node/.style={font=\small}]
			\tikzset{bola/.style={thick, fill=StripA!40}}
			% estable
			\begin{scope}
				\draw[thick, domain=-1.3:1.3, samples=30] plot (\x,{0.9*\x*\x});
				\draw[bola] (0,0.2) circle (0.2);
				\node at (0,-0.6) {(a) Estable};
			\end{scope}
			% indiferente
			\begin{scope}[xshift=4.5cm]
				\draw[thick] (-1.3,0) -- (1.3,0);
				\draw[bola] (0,0.2) circle (0.2);
				\node at (0,-0.6) {(b) Indiferente};
			\end{scope}
			% inestable
			\begin{scope}[xshift=9cm, yshift=0.9cm]
				\draw[thick, domain=-1.3:1.3, samples=30] plot (\x,{-0.9*\x*\x});
				\draw[bola] (0,0.2) circle (0.2);
				\node at (0,-1.5) {(c) Inestable};
			\end{scope}
		\end{tikzpicture}
	\end{center}
	\leyenda{fig:estabilidad}{(a) Tras la perturbación, el centro de gravedad
		sube y el cuerpo \emph{vuelve}. (b) La altura no cambia: se queda en la nueva
		posición. (c) El centro de gravedad baja y el cuerpo \emph{se aleja}.}
	
	% =========================================================
	\section{Procedimiento}
	% =========================================================
	
	Midan antes lo que van a necesitar: la separación $a$ entre orificios, la
	resolución de la regla ($\Delta_{\text{reg}}$), la del dinamómetro
	($\Delta F_{\text{din}}$) y la de la balanza ($\Delta m$). Las distancias se
	miden \textbf{con la regla, entre centros de orificios}, siempre desde el
	apoyo.
	
	\subsection{Parte A: la barra sola}
	
	\begin{enumerate}
		\item Midan la masa $m_b$ de la barra con la balanza y calculen
		$P_b = m_b\,g$.
		\item Cuelguen la barra del dinamómetro por el orificio central y lean su
		peso (Figura~\ref{fig:barra1}). Comparen con $m_b\,g$.
		\item Si la barra queda horizontal colgada de ese orificio, su centro de
		gravedad está allí. Si se inclina, el centro de gravedad está del lado
		hacia el que baja.
	\end{enumerate}
	
	\vspace{0.1cm}
	\begin{center}
		\begin{tikzpicture}[scale=0.62, every node/.style={font=\small}]
			\barra{0}{0}
			\dinaarriba{5.5}{2.6}
			\draw[peso] (5.5,-0.25) -- (5.5,-1.35) node[right] {$P_b$};
			\node[right, align=left, text width=5.2cm] at (11.6,1.0)
			{La lectura del dinamómetro es el peso de la barra};
		\end{tikzpicture}
	\end{center}
	\leyenda{fig:barra1}{La barra colgada del dinamómetro por el orificio central.}
	
	\subsection{Parte B: las tres palancas}
	
	\begin{enumerate}
		\setcounter{enumi}{3}
		\item Armen las tres configuraciones de la Figura~\ref{fig:palancas}.
		Cuelguen en todas la \textbf{misma masa} $m$. El apoyo va en el orificio
		\textbf{central} en la palanca de 1.er género y en un \textbf{extremo} en
		las de 2.º y 3.er género.
		\item En cada una, \textbf{antes de leer el dinamómetro}, midan $d_m$, $d_F$
		y $d_b$, y \emph{predigan} la lectura $F_{\text{pred}}$ con la
		ecuación~(\ref{eq:Fpred}).
		\item Lean el dinamómetro ($F_{\text{med}}$) con la barra horizontal.
	\end{enumerate}
	
	\notabox{La barra \textbf{debe quedar horizontal} cuando lean el
		dinamómetro: si está inclinada, los brazos cambian y la predicción deja de
		valer.}
	
	\begin{center}
		\begin{tikzpicture}[scale=0.62, every node/.style={font=\small}]
			% ---------- (a) 1er genero ----------
			\begin{scope}
				\barra{0}{0}
				\fulcro{5.5}
				\pesa{1.5}{-0.25}{$m$}
				\dinaabajo{9.5}{-2.6}
				\draw[dim] (1.5,0.75) -- (5.5,0.75) node[midway, above, font=\scriptsize] {$d_m$};
				\draw[dim] (5.5,0.75) -- (9.5,0.75) node[midway, above, font=\scriptsize] {$d_F$};
				\node[right, align=left, text width=4.3cm] at (11.6,0)
				{\textbf{(a) 1.er género}\\ apoyo entre la masa\\ y el dinamómetro\\ (el dinamómetro tira\\ hacia abajo)};
			\end{scope}
			% ---------- (b) 2do genero ----------
			\begin{scope}[yshift=-6.0cm]
				\barra{0}{0}
				\fulcro{0.5}
				\pesa{5.5}{-0.25}{$m$}
				\dinaarriba{10.5}{2.2}
				\draw[dim] (0.5,0.75) -- (5.5,0.75) node[midway, above, font=\scriptsize] {$d_m$};
				\draw[dim] (0.5,1.4) -- (10.5,1.4) node[midway, above, font=\scriptsize] {$d_F$};
				\node[right, align=left, text width=4.3cm] at (11.6,0)
				{\textbf{(b) 2.º género}\\ la masa queda entre\\ el apoyo y el\\ dinamómetro};
			\end{scope}
			% ---------- (c) 3er genero ----------
			\begin{scope}[yshift=-11.6cm]
				\barra{0}{0}
				\fulcro{0.5}
				\dinaarriba{3.5}{2.2}
				\pesa{10.5}{-0.25}{$m$}
				\draw[dim] (0.5,1.4) -- (3.5,1.4) node[midway, above, font=\scriptsize] {$d_F$};
				\draw[dim] (0.5,0.75) -- (10.5,0.75) node[midway, above, font=\scriptsize] {$d_m$};
				\node[right, align=left, text width=4.3cm] at (11.6,0)
				{\textbf{(c) 3.er género}\\ el dinamómetro queda\\ entre el apoyo y\\ la masa (como el\\ bíceps en el antebrazo)};
			\end{scope}
		\end{tikzpicture}
	\end{center}
	\leyenda{fig:palancas}{Las tres palancas del laboratorio (las posiciones de los
		orificios son un ejemplo). En (a) el dinamómetro se ancla desde abajo; en
		(b) y (c) cuelga del soporte superior. $d_m$ y $d_F$ son las distancias del
		apoyo a la masa y al dinamómetro. El peso de la barra actúa en su centro de
		gravedad en las tres; $d_b$ es su distancia al apoyo (cero en la (a) si el
		apoyo está en el orificio central).}
	
	\subsection{Parte C: equilibrio estable, indiferente e inestable}
	
	\begin{enumerate}
		\setcounter{enumi}{6}
		\item Pongan la barra en posición \textbf{vertical} y suspéndanla por un
		orificio, sin masas, de modo que pueda girar libremente
		(Figura~\ref{fig:regla}). Pruébenlo con un orificio \emph{arriba} del centro de
		gravedad, con el orificio \emph{del centro} y con uno \emph{abajo} del
		centro. En cada caso, desvíen suavemente la barra y observen qué hace al
		soltarla.
	\end{enumerate}
	
	\vspace{0.1cm}
	\begin{center}
		\begin{tikzpicture}[scale=0.95, every node/.style={font=\small}]
			\newcommand{\reglav}[3]{%
				% #1 = x del panel, #2 = coordenada del apoyo en la barra, #3 = etiqueta
				\draw[line width=3pt, gray!60] (#1-1.0,0) -- (#1,0);
				\begin{scope}[shift={(#1,{-0.42*#2})}, rotate=90, scale=0.42]
					\barra{0}{0}
				\end{scope}
				\draw[thick, fill=white] (#1,0) circle (0.09);
				\fill[StripA] (#1,{0.42*(5.5-#2)}) circle (0.07);
				\draw[peso] (#1+0.3,{0.42*(5.5-#2)}) -- (#1+0.3,{0.42*(5.5-#2)-0.9}) node[right] {$P_b$};
				\draw[->, thick, StripB] (#1-1.0,{-0.42*#2+0.15}) -- (#1-0.35,{-0.42*#2+0.15});
				\node[align=center] at (#1,-4.35) {#3};
			}
			\reglav{0}{8.5}{(A) apoyo\\ sobre el CG}
			\reglav{4.4}{5.5}{(B) apoyo\\ en el CG}
			\reglav{8.8}{2.5}{(C) apoyo\\ bajo el CG}
		\end{tikzpicture}
	\end{center}
	\leyenda{fig:regla}{Barra vertical suspendida de un perno (círculo blanco) en
		tres posiciones. El punto rojo es el centro de gravedad (CG); la flecha azul
		es la perturbación. ¿Qué tipo de equilibrio es cada caso? Es una de las
		preguntas del análisis.}
	
	\noindent
	En el caso (C) el equilibrio es muy delicado: sostengan la barra con la mano
	hasta soltarla despacio.
	
	\vspace{0.2cm}
	\begin{center}
		\small
		$m_b = $ \rule{1.3cm}{0.4pt} kg
		\quad
		$P_b = $ \rule{1.3cm}{0.4pt} N
		\quad
		$a = $ \rule{1.3cm}{0.4pt} cm
		\quad
		$m = $ \rule{1.3cm}{0.4pt} kg
		
		\vspace{0.15cm}
		$\Delta_{\text{reg}} = $ \rule{1.3cm}{0.4pt} cm
		\qquad
		$\Delta F_{\text{din}} = $ \rule{1.3cm}{0.4pt} N
		\qquad
		$\Delta m = $ \rule{1.3cm}{0.4pt} kg
		
		\vspace{0.4cm}
		\setlength{\tabcolsep}{4pt}
		\begin{tabular}{|l|c|c|c|c|c|}
			\hline
			Montaje & $d_m$ (cm) & $d_F$ (cm) & $d_b$ (cm)
			& $F_{\text{pred}}$ (N) & $F_{\text{med}}$ (N) \\
			\hline
			\rule{0pt}{3ex} 1.er género & & & & & \\ \hline
			\rule{0pt}{3ex} 2.º género & & & & & \\ \hline
			\rule{0pt}{3ex} 3.er género & & & & & \\ \hline
		\end{tabular}
	\end{center}
	
	% =========================================================
	\section{Análisis de resultados (en casa)}\label{sec:analisis}
	% =========================================================
	
	\subsection{Predicción contra medida}
	
	Para cada montaje, calculen la discrepancia:
	\[
	\text{disc.} = \frac{|F_{\text{med}} - F_{\text{pred}}|}{F_{\text{pred}}}\times 100\%
	\]
	
	\noindent
	Calculen también la incertidumbre de $F_{\text{pred}}$ siguiendo las reglas
	del Taller~2, en este orden:
	\begin{enumerate}
		\item Cada distancia es la resta de dos lecturas de la regla:
		$\Delta d = \sqrt{2}\,\Delta_{\text{reg}}$. Y $\Delta P = g\,\Delta m$.
		\item Cada torque es un \emph{producto} ($P_m\,d_m$ y $P_b\,d_b$): sumen
		en cuadratura las incertidumbres \emph{relativas} y multipliquen por el
		valor del producto.
		\item El numerador es una \emph{suma}: sumen en cuadratura las
		incertidumbres \emph{absolutas} de los dos torques.
		\item Al dividir entre $d_F$ (un \emph{cociente}): sumen en cuadratura la
		incertidumbre relativa del numerador y la de $d_F$.
	\end{enumerate}
	
	\noindent
	Digan si $F_{\text{pred}}$ y $F_{\text{med}}$ son compatibles: es decir, si
	su diferencia es menor que la raíz de la suma de los cuadrados de sus dos
	incertidumbres.
	
	\subsection{Preguntas}
	
	\begin{enumerate}
		\item ¿A qué clase de equilibrio corresponde la barra en cada posición de
		la Figura~\ref{fig:regla}? Justifiquen con lo que pasa con la altura del
		centro de gravedad.
		\item ¿El centro de gravedad de la barra coincide con el orificio central?
		¿Cómo lo comprobaron en la Parte A?
		\item Si la lectura del dinamómetro fuera desconocida, ¿podría
		calcularse? ¿Con qué ecuación?
		\item ¿En qué casos es importante incluir el peso de la barra? ¿En qué
		casos podría despreciarse? Usen sus datos: calculen cuánto cambia
		$F_{\text{pred}}$ si se omite el término $P_b\,d_b$.
		\item En su palanca de 3.er género, ¿cuántas veces es mayor la fuerza del
		dinamómetro que el peso de la masa? Un bíceps se inserta a unos 4~cm del
		codo y la mano queda a unos 35~cm (valores aproximados). ¿Cuántas veces el
		peso que sostiene debe hacer fuerza el bíceps?
		\item Clasifiquen como palancas los sistemas de la tabla siguiente.
	\end{enumerate}
	
	\begin{center}
		\small
		\begin{tabular}{|l|c|c|c|c|}
			\hline
			Sistema & Apoyo & Carga & Esfuerzo & Género \\
			\hline
			\rule{0pt}{3ex} Balanza de platillos & & & & \\ \hline
			\rule{0pt}{3ex} Tijeras & & & & \\ \hline
			\rule{0pt}{3ex} Brazo (flexión del codo) & & & & \\ \hline
			\rule{0pt}{3ex} Carretilla & & & & \\ \hline
			\rule{0pt}{3ex} Mordida (mandíbula) & & & & \\ \hline
		\end{tabular}
	\end{center}
	
	\noindent
	Para el brazo: el apoyo es el codo, la carga es el peso de la mano y lo que
	sostiene, y el esfuerzo es la fuerza del bíceps, que se inserta en el
	antebrazo muy cerca del codo.
	
	% =========================================================
	\section{Conclusiones}
	% =========================================================
	
	Presenten conclusiones contrastando los resultados experimentales con las
	condiciones de equilibrio. Además de sus conclusiones propias, discutan:
	
	\begin{itemize}
		\item ¿Se cumplieron las condiciones de equilibrio dentro de las
		incertidumbres?
		\item ¿Qué montaje tuvo mayor discrepancia? ¿Qué lo explicaría (barra
		inclinada, lectura del dinamómetro, posición de los orificios, centro
		de gravedad que no está en el centro)?
		\item ¿Qué implicación tiene, para los músculos y las articulaciones, que
		el antebrazo sea una palanca de 3.er género?
	\end{itemize}
	
	% =========================================================
	\section{Formato del Informe 1}\label{sec:informe}
	% =========================================================
	
	Esta práctica se entrega como \textbf{informe} (PDF), no como bitácora.
	Una sola entrega por grupo, con el nombre
	\texttt{Informe1\_GrupoN.pdf}. Debe tener:
	
	\begin{enumerate}
		\item \textbf{Introducción y marco teórico:} equilibrio, torque y
		palancas, con las ecuaciones que usarán. Incluyan por qué importa en
		biomecánica.
		\item \textbf{Metodología:} montaje con un esquema propio, instrumentos y
		sus incertidumbres.
		\item \textbf{Resultados:} tablas con unidades e incertidumbres, y el
		cálculo de $F_{\text{pred}}$ de cada montaje.
		\item \textbf{Análisis y discusión:} la comparación y las preguntas de la
		Sección~\ref{sec:analisis}.
		\item \textbf{Conclusiones.}
		\item \textbf{Referencias} en un formato consistente (al menos dos).
	\end{enumerate}
	
	\notabox{La primera versión se entrega el \textbf{3 de noviembre de 2026}.
		Recibirán comentarios y entregarán una versión corregida; las fechas
		exactas se publicarán en Classroom.}
	
	% =========================================================
	\section{Referencias}
	% =========================================================
	
	\begin{itemize}
		\item L. Agulles Perdós, \emph{Guías de Laboratorio: Mecánica y Ondas
			para Biociencias}, Departamento de Física, Universidad Nacional de
		Colombia, 2026. Laboratorio 9.
	\end{itemize}
	
	\vfill
	\noindent
	{\footnotesize \textit{Auxiliar docente: Miguel Andrés Perdomo Gutiérrez
			--- mperdomog@unal.edu.co}}
	
\end{document}